% Shared authored contribution: included by the active WP1 chapter and legacy fiche.
\paragraph{CGAL mesh smoothing and CGALSCALE (2026).}
François Protais (Inria, TITANE) developed volumetric mesh smoothing
methods that are particularly effective for handling boundary elements.
The approach addresses needs within CGAL's volume-mesh workflows.
The resulting \texttt{Mesh\_smoothing\_3} package optimizes vertex positions
while retaining mesh connectivity. It has been integrated into CGAL's
\texttt{main} branch and is scheduled for CGAL~6.3. The
\href{https://cgal.geometryfactory.com/CGAL/doc/main/Mesh_smoothing_3/index.html#Chapter_3D_Mesh_smoothing}{package documentation}
provides the interface and examples; the
\href{https://www.cgal.org/2026/10/01/mesh_smoothing_3/}{official CGAL announcement}
was published on 1 October 2026.

CGALSCALE (CGAL Components at All Scales) is a new i-Lab partnership
between the Inria TITANE team and GeometryFactory. The project aims to
adapt CGAL's geometric algorithms to emerging computing environments by
developing distributed-memory and GPU-enabled versions of key CGAL
components. Its objective is twofold: to process increasingly large
geometric datasets and to exploit computing platforms ranging from
low-power embedded devices to hybrid CPU--GPU systems, cloud
infrastructures and supercomputers. A central challenge is to preserve
CGAL's robustness and exact geometric computation paradigm while
achieving scalability across data size and computational resources.
The partnership and its development objectives are described in the
\href{https://jobs.inria.fr/public/classic/classic/offres/2026-10293}{Inria CGALSCALE project description}.
